\documentclass[11pt, a4paper]{article}

% --- PAQUETES PRINCIPALES ---
\usepackage[utf8]{inputenc}
\usepackage[spanish, es-tabla, es-noquoting]{babel} % Vital para evitar conflictos TikZ/Babel
\usepackage[margin=2.5cm]{geometry}
\usepackage{amsmath, amssymb, mathtools}
\usepackage{siunitx}
\usepackage{graphicx}
\usepackage{booktabs}
\usepackage{tabularx}
\usepackage[most]{tcolorbox}
\usepackage{xcolor}

% --- PAQUETES PARA CIRCUITOS Y GRÁFICOS ---
\usepackage[american]{circuitikz}
\usetikzlibrary{shapes.geometric, arrows, positioning, babel}

% --- CONFIGURACIÓN DE SIUNITX ---
\sisetup{
    output-decimal-marker = {,},
    per-mode = symbol,
    inter-unit-product = \ensuremath{{}\cdot{}}
}

% --- CABECERAS Y PIES DE PÁGINA ---
\usepackage{fancyhdr}
\setlength{\headheight}{14.5pt} % Evita warning de overfull header
\pagestyle{fancy}
\fancyhf{}
\lhead{\textbf{Circuitos Electrónicos III (IMT-221)}}
\rhead{Anexo Técnico: BOM y Protocolo Shunt}
\cfoot{\thepage}
\renewcommand{\headrulewidth}{0.4pt}

% --- ENTORNOS PERSONALIZADOS ---
\newtcolorbox{aviso}[1][]{
  colback=red!5!white,
  colframe=red!75!black,
  fonttitle=\bfseries,
  title=#1
}
\newtcolorbox{cajaopciones}[1][]{
  colback=blue!3!white,
  colframe=blue!80!black,
  fonttitle=\bfseries,
  title=#1
}
\newtcolorbox{cajaevaluacion}[1][]{
  colback=gray!10!white,
  colframe=black!75,
  fonttitle=\bfseries,
  title=#1
}

\begin{document}

\begin{center}
    {\LARGE \textbf{ANEXO TÉCNICO AL HITO 1}}\\[0.2cm]
    {\Large \textbf{LISTA DE MATERIALES (BOM) ACTUALIZADA Y PROTOCOLO DE CONTINGENCIA PARA RESISTENCIAS DE SENSADO ($R_{\text{shunt}}$)}}\\[0.4cm]
    
    UNIVERSIDAD CATÓLICA BOLIVIANA ``SAN PABLO'' -- SEDE TARIJA\\[0.1cm]
    Departamento de Ciencias Exactas e Ingeniería\\[0.1cm]
    Carrera: Ingeniería Mecatrónica / Ingeniería Biomédica\\[0.1cm]
    Asignatura: Circuitos Electrónicos III (IMT-221) | Gestión: Semestre 2-2026\\[0.1cm]
    \textit{Docente: M.Sc. Ing. Bernardo Quiroga Turdera}
    \vspace{0.3cm}
    \hrule
\end{center}

\vspace{0.4cm}

% ==============================================================================
% 1. LISTA DE MATERIALES
% ==============================================================================
\section{Lista de Materiales Oficial y Consolidada (BOM Hito 1)}

La siguiente tabla detalla los componentes requeridos para la implementación, simulación y validación experimental de la etapa de protección y supresión de transitorios inductivos del Hito 1:

\begin{table}[h!]
\centering
\caption{Especificaciones de Componentes para el Hito 1 del PFI}
\label{tab:bom_hito1}
\footnotesize
\renewcommand{\arraystretch}{1.4}
% Ajuste de anchos para evitar Overfull \hbox
\begin{tabularx}{\textwidth}{@{} c >{\bfseries\raggedright\arraybackslash}p{3cm} >{\raggedright\arraybackslash}p{4.5cm} c X @{}}
\toprule
\textbf{Ítem} & \textbf{Componente / Denominación} & \textbf{Parámetros Críticos y Tolerancia} & \textbf{Cant.} & \textbf{Función Técnica en el Circuito} \\
\midrule
1 & Diodo Zener de Precisión & BZX55C5V1 ($V_Z = 5.1\text{ V} \pm 5\%$, $P_{ZM} = 500\text{ mW}$, $r_z \le 15\,\Omega$) & 1 & Limitador de techo de tensión segura para ADC ($0 - 5.1\text{ V}$). \\
2 & Diodo Schottky de Señal & BAT42, BAT43 o SD101 ($V_F \le 0.33\text{ V}$, $t_{rr} < 1\text{ ns}$) & 1 & Alternativa Clamping Activo (Rail-to-Rail) al riel de $+3.3\text{ V}$. \\
3 & Diodos Rápidos de Conmutación & 1N4148 ($t_{rr} \le 4\text{ ns}$, $I_{F(\max)} = 200\text{ mA}$, $V_{R} = 100\text{ V}$) & 3 a 5 & Fijación de piso a $-0.7\text{ V}$, redundancia o cadena de diodos (string). \\
4 & Diodo Rectificador de Potencia & 1N4007 ($I_{F(\text{avg})} = 1\text{ A}$, $I_{FSM} = 30\text{ A}$, $PIV = 1000\text{ V}$) & 1 & Diodo Flyback de supresión de energía inductiva en la carga motriz. \\
5 & Resistor Limitador de Serie ($R_s$) & $180\,\Omega$ o $220\,\Omega$ / $0.5\text{ W}$ a $1\text{ W}$ (Película metálica $\pm 5\%$) & 1 & Absorción de sobretensión y limitación de corriente hacia el Zener. \\
6 & Resistor de Sensado Shunt ($R_{\text{shunt}}$) & $0.1\,\Omega$ / $2\text{ W}$ a $5\text{ W}$ (Cerámica/cemento, tolerancia $\le 5\%$) & 1 & Elemento transductor corriente-voltaje ($100\text{ mV/A}$). \\
7 & Resistores Auxiliares / Cargas & $10\,\Omega, 100\,\Omega, 1\text{ k}\Omega, 10\text{ k}\Omega, 47\text{ k}\Omega$ ($1/4\text{ W}$) & Kit & Resistencia de carga de prueba ($R_{LOAD}$) y divisores de inyección. \\
8 & Actuador Inductivo de Prueba & Motorreductor DC "Amarillo" ($3\text{ V} - 12\text{ V}$) o Solenoide inductivo & 1 & Generador físico de transitorio inductivo por conmutación. \\
\bottomrule
\end{tabularx}
\end{table}

% ==============================================================================
% 2. EL RESISTOR SHUNT
% ==============================================================================
\section{El Resistor Shunt ($R_{\text{shunt}}$): Rol en el PFI y Requisito Nominal}

En la arquitectura del Proyecto Final Integrador, el resistor Shunt ($R_{\text{shunt}}$) es el elemento primario encargado de traducir la corriente consumida por la carga en una diferencia de potencial analógica:
\begin{equation}
    V_{\text{shunt}}(t) = I_{\text{carga}}(t) \cdot R_{\text{shunt}}
\end{equation}

\textbf{Especificación de Diseño para el Sistema Integrado:}
\begin{itemize}
    \item \textbf{Valor Nominal Óptimo:} $R_{\text{shunt}} = 0.10\,\Omega$ ($100\text{ m}\Omega$).
    \item \textbf{Caída de Tensión:} A corriente nominal de $1.0\text{ A}$, genera una señal de $100\text{ mV}$, preservando el $99.2\%$ de la tensión de la fuente de $+12\text{ V}$ para el actuador.
    \item \textbf{Disipación Térmica:} $P = I^2 R = (1.5\text{ A})^2 \cdot 0.1\,\Omega = \mathbf{0.225\text{ W}}$ (un resistor de $5\text{ W}$ opera con un factor de seguridad térmico superior al $95\%$).
\end{itemize}

\newpage
% ==============================================================================
% 3. PROTOCOLO DE CONTINGENCIA
% ==============================================================================
\section{Protocolo de Contingencia: Alternativas ante la ausencia de $R_{\text{shunt}} = 0.1\,\Omega$}

Si el equipo de trabajo no dispone de la resistencia de cemento de $0.1\,\Omega$ / $5\text{ W}$ durante la sesión del Hito 1, deberá aplicar obligatoriamente una de las siguientes cuatro alternativas de ingeniería para validar la red de protección sin alterar los objetivos de aprendizaje:

\vspace{0.3cm}
\begin{cajaopciones}[Alternativas de Contingencia para Banco de Trabajo]
    \begin{itemize}
        \item[\textbf{A.}] \textbf{Síntesis en Paralelo} $\longrightarrow$ Conectar 5 a 10 resistores estándar de 1/4W en paralelo.
        \item[\textbf{B.}] \textbf{Inyección Pura de Señal} $\longrightarrow$ Reemplazar $R_{\text{shunt}}$ por $10\,\Omega$, $100\,\Omega$ o $1\text{ k}\Omega$ (Sin motor).
        \item[\textbf{C.}] \textbf{Supresión Flyback Directa} $\longrightarrow$ Probar Flyback con motor a 12V sin resistor en serie.
        \item[\textbf{D.}] \textbf{Shunt Calibrado de Cobre} $\longrightarrow$ Construir shunt resistivo con alambre de cobre esmaltado.
    \end{itemize}
\end{cajaopciones}
\vspace{0.3cm}

\subsection*{Alternativa A: Síntesis de Shunt por Arreglo en Paralelo (Recomendada para Potencia)}
Si se dispone de resistencias comunes de película de carbón o metálica de $1/4\text{ W}$ o $1/2\text{ W}$, se puede construir una resistencia equivalente de bajo valor óhmico y mayor capacidad de disipación de potencia conectando $N$ resistores idénticos en paralelo:
\begin{equation}
    R_{eq} = \frac{R}{N} \qquad P_{total} = N \cdot P_{individual}
\end{equation}

\begin{itemize}
    \item \textbf{Opción A.1:} $10$ resistencias de $1.0\,\Omega$ / $1/4\text{ W}$ en paralelo:
    $$R_{eq} = \frac{1.0\,\Omega}{10} = \mathbf{0.10\,\Omega} \qquad P_{total} = 10 \times 0.25\text{ W} = \mathbf{2.50\text{ W}}$$
    
    \item \textbf{Opción A.2:} $5$ resistencias de $0.47\,\Omega$ / $0.5\text{ W}$ en paralelo:
    $$R_{eq} = \frac{0.47\,\Omega}{5} = \mathbf{0.094\,\Omega} \qquad P_{total} = 5 \times 0.50\text{ W} = \mathbf{2.50\text{ W}}$$
    
    \item \textbf{Opción A.3:} $5$ resistencias de $1.0\,\Omega$ / $1/4\text{ W}$ en paralelo:
    $$R_{eq} = \frac{1.0\,\Omega}{5} = \mathbf{0.20\,\Omega} \qquad P_{total} = 5 \times 0.25\text{ W} = \mathbf{1.25\text{ W}}$$
\end{itemize}

\subsection*{Alternativa B: Emulación por Divisor Resistivo (Validación de Protección de Entrada)}
El objetivo del Hito 1 es evaluar la red de sujeción (clamping) y limitación (clipping) ante sobretensiones críticas. Si no se opera la planta motriz en esta etapa, el esquema de prueba será:

\vspace{0.3cm}
\begin{figure}[h!]
    \centering
    \begin{circuitikz}[american, scale=0.9, transform shape]
        % Generador
        \draw (0,0) to[sV, l={Generador de Señales}] (0,-3) node[ground]{};
        \node[above, text width=4cm, align=center] at (0, 0.5) {\footnotesize $15\text{ V}_{pp}$, $\qty{1}{\kilo\hertz}$, $+2.5\text{V offset}$};
        
        % Divisor (Uso llaves {} para proteger los textos)
        \draw (0,0) -- (1.5,0) to[R, l_={$R_{DIV}$ ($1\text{k}\Omega$ o $10\text{k}\Omega$)}] (4.5,0);
        
        % Nodo Central
        \draw (4.5,0) to[short, -*] (5.5,0);
        \node[above] at (5.5, 0.1) {\textbf{Nodo de Sensado ($v_{in}$)}};
        
        % Resistencia simulada a GND (¡LLAVES PROTECTORAS AQUÍ PARA LA COMA!)
        \draw (5.5,0) to[R, l={$R_{\text{shunt\_sim}}$ ($10\,\Omega$, $100\,\Omega$ o $1\text{k}\Omega$)}] (5.5,-3) node[ground]{};
        
        % A la red
        \draw (5.5,0) -- (7.5,0) to[R, l=$R_s$] (10,0) to[short, -o] (11,0) node[right]{\textbf{[ Red Diodos ]}};
    \end{circuitikz}
    \caption{Emulación de señal mediante divisor resistivo.}
\end{figure}
\vspace{0.2cm}

\textbf{Procedimiento:} Se sustituye $R_{\text{shunt}}$ por una resistencia cualquiera del kit ($10\,\Omega, 100\,\Omega \text{ o } 1\text{ k}\Omega$).\\
\textbf{Validación:} El generador de funciones inyecta directamente la sobretensión de prueba ($15\text{ V}_{pp}$). Se verifica en el osciloscopio que la tensión entregada al nodo del ADC se mantiene dentro de la cota segura ($0.0\text{ V} \le V_{ADC} \le 5.1\text{ V}$ o $3.3\text{ V}$).

\newpage
\subsection*{Alternativa C: Ensayo Desacoplado del Diodo Flyback}
Si se desea verificar la supresión de la fuerza contraelectromotriz ($V_L = L \frac{di}{dt}$) del motor DC real:

\begin{aviso}[Advertencia de Seguridad en Banco]
NO conectar resistencias de $1/4\text{ W}$ de $100\,\Omega$ o $1\text{ k}\Omega$ en serie con el motor conectado a $+12\text{ V}$. Riesgo inminente de destrucción por disipación térmica y quemadura de protoboard.
\end{aviso}

\begin{figure}[h!]
    \centering
    \begin{circuitikz}[american, scale=0.85, transform shape]
        \draw (0,3) node[vcc]{+12V} -- (0,2);
        \draw (-1.5,2) -- (1.5,2);
        \draw (-1.5,2) to[generic, l={Motor DC}] (-1.5,0);
        \draw (1.5,2) to[D, l={1N4007 (Flyback)}, invert] (1.5,0);
        \draw (-1.5,0) -- (1.5,0);
        \draw (0,0) to[cute open switch, l={Switch / Transistor}] (0,-2) node[ground]{};
    \end{circuitikz}
    \caption{Configuración segura para ensayo de Diodo Flyback directo.}
\end{figure}

\textbf{Procedimiento Seguro:}
\begin{enumerate}
    \item Conectar el motor directamente entre la fuente de $+12\text{ V}$ y el switch mecánico hacia GND.
    \item Conectar el diodo 1N4007 en antiparalelo directo a los bornes del motor (cátodo a $+12\text{ V}$, ánodo al nodo de conmutación).
    \item Medir con el Canal 1 del osciloscopio el pico de apertura con y sin el diodo 1N4007 para evidenciar la absorción del transitorio inductivo.
\end{enumerate}

\subsection*{Alternativa D: Calibración de Shunt por Longitud de Alambre de Cobre}
En aplicaciones de instrumentación donde no se cuenta con componentes comerciales, es técnicamente válido confeccionar un resistor shunt calibrado utilizando un segmento de alambre de cobre esmaltado (bobinado) aprovechando su resistividad:
\begin{equation}
    R = \rho \cdot \frac{L}{A}
\end{equation}
Para alambre de cobre magneto calibre AWG 24 (sección $A \approx 0.205\text{ mm}^2$, resistividad $\rho \approx 1.72 \times 10^{-8}\,\Omega\cdot\text{m}$):
\begin{equation*}
    \text{Resistencia específica} \approx 0.084\,\Omega/\text{metro}
\end{equation*}
\textbf{Construcción:} Un segmento de $1.19\text{ metros}$ de alambre AWG 24 enrollado en forma no inductiva (bifilar) proporciona exactamente $0.10\,\Omega$ con alta capacidad de corriente.

% ==============================================================================
% 4. MATRIZ DE PRECAUCIONES
% ==============================================================================
\section{Matriz de Precauciones en Banco de Trabajo}

\begin{table}[h!]
\centering
\footnotesize
\renewcommand{\arraystretch}{1.5}
\begin{tabularx}{\textwidth}{@{} >{\raggedright\arraybackslash}X >{\raggedright\arraybackslash}X >{\raggedright\arraybackslash}X @{}}
\toprule
\textbf{Práctica Incorrecta / Riesgo en Banco} & \textbf{Consecuencia Física} & \textbf{Acción Correctiva de Ingeniería} \\
\midrule
Conectar un resistor de $1/4\text{ W}$ ($10\,\Omega$ a $100\,\Omega$) en serie con el motor a $+12\text{ V}$. & 
Destrucción térmica instantánea del resistor por sobrecarga de potencia ($P > 1.4\text{ W}$). & 
Usar arreglo en paralelo (Alternativa A) o conectar el motor sin shunt para prueba de Flyback (Alternativa C). \\

Omitir la unión de tierras entre la fuente de $+12\text{ V}$, el generador y el osciloscopio. & 
Niveles de referencia flotantes, mediciones de offset erróneas y lazos de masa (\textit{ground loops}). & 
Conectar todos los terminales de masa al nodo central de GND Común de la tarjeta de prototipos. \\

Colocar el diodo Zener en polarización directa en lugar de inversa. & 
La salida quedará recortada a $+0.7\text{ V}$ en lugar de regular a $+5.1\text{ V}$. & 
Verificar cátodo (banda negra) conectado al nodo positivo de la señal. \\

Conectar el diodo Schottky BAT43 al revés en el riel de $+3.3\text{ V}$. & 
Cortocircuito permanente del riel de alimentación del microcontrolador a través del diodo en directa. & 
El ánodo del Schottky debe ir al nodo analógico del ADC y el cátodo al riel de $+3.3\text{ V}$. \\
\bottomrule
\end{tabularx}
\end{table}

\vspace{0.4cm}
\begin{cajaevaluacion}[Nota de Evaluación para el Hito 1]
Los equipos que utilicen cualquiera de las alternativas de contingencia en la sesión práctica deberán documentar explícitamente en el \textbf{Pilar 1} de su memoria técnica (Jupyter Notebook) la justificación analítica del valor equivalente empleado y la deducción de potencia disipada. Para el Hito 2 (Sensado por Par Diferencial BJT), la incorporación del resistor nominal de $0.1\,\Omega$ / $5\text{ W}$ será obligatoria.
\end{cajaevaluacion}

\end{document}
